% Part: first-order-logic
% Chapter: syntax-and-semantics
% Section: Semantic Notions

\documentclass[../../../include/open-logic-section]{subfiles}

\begin{document}

\olfileid{fol}{syn}{sem}
\olsection{المفاهيم الدلالية}

\begin{explain}
بعد تعريف ال!!{structure}s للغات الرتبة الأولى، يمكننا تعريف بعض
الخصائص الدلالية الأساسية للجمل وبعض العلاقات بينها. وأبسطها مفهوم
\emph{صلاحية} الجملة. تكون الجملة صالحة إذا كانت مُرضاة في كل
!!{structure}. فالجمل الصالحة هي التي تُرضى كيفما أُوّلت الرموز غير
المنطقية فيها. ولذلك تسمى الجمل الصالحة أيضًا \emph{حقائق
منطقية}---فهي صادقة، أي مُرضاة، في أي !!{structure}، ومن ثم لا يعتمد
صدقها إلا على الرموز المنطقية الواردة فيها وال!!{structure} التركيبية لها، لا على
الرموز غير المنطقية أو تأويلها.
\end{explain}

\begin{defn}[الصلاحية]
تكون الجملة~$!A$ \emph{صالحة}، $\Entails !A$، إذا وفقط إذا كان
$\Sat{M}{!A}$ لكل !!{structure}~$\Struct M$.
\end{defn}

\begin{defn}[الاستلزام]
\emph{تستلزم} مجموعة من الجمل~$\Gamma$ جملة~$!A$،
$\Gamma \Entails !A$، إذا وفقط إذا كان، لكل
!!{structure}~$\Struct M$ تحقق $\Sat{M}{\Gamma}$،
$\Sat{M}{!A}$.
\end{defn}


\begin{defn}[قابلية الإرضاء]
تكون مجموعة من الجمل~$\Gamma$ \emph{قابلة للإرضاء} إذا كان
$\Sat{M}{\Gamma}$ ل!!{structure} ما~$\Struct M$. وإذا لم تكن
$\Gamma$ قابلة للإرضاء سميت \emph{غير قابلة للإرضاء}.
\end{defn}

\begin{prop}
تكون الجملة~$!A$ صالحة إذا وفقط إذا كان
$\Gamma \Entails !A$ لكل مجموعة من الجمل~$\Gamma$.
\end{prop}

\begin{proof}
في الاتجاه الأمامي، لتكن $!A$ صالحة، ولتكن $\Gamma$ مجموعة من
الجمل. ولتكن $\Struct M$ !!a{structure} بحيث
$\Sat{M}{\Gamma}$. لما كانت $!A$ صالحة، كان $\Sat{M}{!A}$، ومن
ثم $\Gamma \Entails !A$.

ولإثبات مقابل عكس الاتجاه العكسي، لتكن $!A$ غير صالحة، بحيث توجد
!!a{structure}~$\Struct M$ تحقق $\Sat/{M}{!A}$. عندما
$\Gamma = \{ \ltrue \}$، يكون $\Sat{M}{\Gamma}$ لأن
$\ltrue$ صالحة. ومن ثم توجد !!a{structure}~$\Struct M$ بحيث
$\Sat{M}{\Gamma}$ لكن $\Sat/{M}{!A}$، ولذلك لا تستلزم~$\Gamma$
الصيغة~$!A$.
\end{proof}

\begin{prop}
\ollabel{prop:entails-unsat}
$\Gamma \Entails !A$ إذا وفقط إذا كانت
$\Gamma \cup \{\lnot !A\}$ غير قابلة للإرضاء.
\end{prop}

\begin{proof}
في الاتجاه الأمامي، افترض أن $\Gamma \Entails !A$، وافترض، خلافًا
للمطلوب، أن ثمة !!a{structure}~$\Struct M$ بحيث
$\Sat{M}{\Gamma \cup \{ \lnot !A \}}$. لما كان
$\Sat{M}{\Gamma}$ و$\Gamma \Entails !A$، كان $\Sat{M}{!A}$.
وكذلك، لما كان $\Sat{M}{\Gamma\cup \{ \lnot !A \}}$، كان
$\Sat{M}{\lnot !A}$، فيكون لدينا كل من $\Sat{M}{!A}$ و
$\Sat/{M}{!A}$، وهذا تناقض. ومن ثم لا يمكن أن توجد
!!{structure}~$\Struct M$ كهذه، فتكون
$\Gamma \cup \{ \lnot !A \}$ غير قابلة للإرضاء.

وفي الاتجاه العكسي، افترض أن $\Gamma \cup \{ \lnot !A \}$ غير
قابلة للإرضاء. إذن، لكل !!{structure}~$\Struct M$، إما
$\Sat/{M}{\Gamma}$ أو $\Sat{M}{!A}$. ومن ثم، لكل !!{structure}
~$\Struct M$ تحقق $\Sat{M}{\Gamma}$، يكون $\Sat{M}{!A}$،
ولذلك $\Gamma \Entails !A$.
\end{proof}

\begin{prob}
\begin{enumerate}
\item بيّن أن $\Gamma \Entails \lfalse$ إذا وفقط إذا كانت $\Gamma$
  غير قابلة للإرضاء.
\item بيّن أن $\Gamma \cup \{!A\} \Entails \lfalse$ إذا وفقط إذا
  كان $\Gamma \Entails \lnot !A$.
\item افترض أن $c$ لا يرد في~$!A$ ولا في~$\Gamma$. بيّن أن
  $\Gamma \Entails \lforall[x][!A]$ إذا وفقط إذا كان
  $\Gamma \Entails \Subst{!A}{c}{x}$.
\end{enumerate}
\end{prob}

\begin{prop}
إذا كان $\Gamma \subseteq \Gamma'$ و$\Gamma \Entails !A$، فإن
$\Gamma' \Entails !A$.
\end{prop}

\begin{proof}
افترض أن $\Gamma \subseteq \Gamma'$ وأن $\Gamma \Entails !A$.
ولتكن $\Struct M$ بنية بحيث $\Sat{M}{\Gamma'}$؛ عندئذ
$\Sat{M}{\Gamma}$، ولما كان $\Gamma \Entails !A$، نحصل على
$\Sat{M}{!A}$. ومن ثم، متى كان $\Sat{M}{\Gamma'}$، كان
$\Sat{M}{!A}$، ولذلك $\Gamma' \Entails !A$.
\end{proof}


\begin{thm}[مبرهنة الاستنباط الدلالية]
\ollabel{thm:sem-deduction}
$\Gamma \cup \{!A\} \Entails !B$ إذا وفقط إذا كان
$\Gamma \Entails !A \lif !B$.
\end{thm}

\begin{proof}
في الاتجاه الأمامي، ليكن
$\Gamma \cup \{ !A \} \Entails !B$، ولتكن $\Struct M$
!!a{structure} بحيث $\Sat{M}{\Gamma}$. إذا كان $\Sat{M}{!A}$،
فإن $\Sat{M}{\Gamma \cup \{ !A \} }$، ولما كانت
$\Gamma \cup \{ !A \}$ تستلزم~$!B$، نحصل على $\Sat{M}{!B}$.
ولذلك $\Sat{M}{!A \lif !B}$، ومن ثم
$\Gamma \Entails !A \lif !B$.

وفي الاتجاه العكسي، ليكن $\Gamma \Entails !A \lif !B$، ولتكن
$\Struct M$ !!a{structure} بحيث
$\Sat{M}{\Gamma \cup \{ !A \}}$. عندئذ $\Sat{M}{\Gamma}$،
فيكون $\Sat{M}{!A \lif !B}$، ولما كان $\Sat{M}{!A}$، يكون
$\Sat{M}{!B}$. ومن ثم، متى كان
$\Sat{M}{\Gamma \cup \{ !A \} }$، كان $\Sat{M}{!B}$، ولذلك
$\Gamma \cup \{ !A \} \Entails !B$.
\end{proof}

\begin{prop}\ollabel{prop:quant-terms}
  لتكن $\Struct{M}$ !!a{structure}، ولتكن $!A(x)$ !!a{formula} ذات
  متغير حر واحد~$x$، وليكن $t$~حدًا مغلقًا. عندئذ:
  \begin{tagenumerate}{prvEx,prvAll}
    \tagitem{prvEx}{$!A(t) \Entails \lexists[x][!A(x)]$}{}
    \tagitem{prvAll}{$\lforall[x][!A(x)] \Entails !A(t)$}{}
  \end{tagenumerate}
\end{prop}

\begin{proof}
  \begin{tagenumerate}{prvEx,prvAll}
    \tagitem{prvEx}{%
      \iftag{probEx}{تمرين.}{افترض $\Sat{M}{!A(t)}$. وليكن $s$
        إسناد متغيرات يحقق $s(x) = \Value{t}{M}$. عندئذ
        $\Sat{M}{!A(t)}[s]$ لأن $!A(t)$ !!a{sentence}. وبحسب
        \olref[ext]{prop:ext-formulas}، يكون
        $\Sat{M}{!A(x)}[s]$. وبحسب
        \olref[ass]{prop:sat-quant}،
        $\Sat{M}{\lexists[x][!A(x)]}$.}}{}
    \tagitem{prvAll}{%
      \iftag{probAll}{تمرين.}{افترض
        $\Sat{M}{\lforall[x][!A(x)]}$. وليكن $s$ إسناد متغيرات
        يحقق $s(x) = \Value{t}{M}$. بحسب
        \olref[ass]{prop:sat-quant}، يكون
        $\Sat{M}{!A(x)}[s]$. وبحسب
        \olref[ext]{prop:ext-formulas}، يكون
        $\Sat{M}{!A(t)}[s]$. وبحسب
        \olref[ass]{prop:sentence-sat-true}، يكون
        $\Sat{M}{!A(t)}$ لأن $!A(t)$ !!a{sentence}.}}{}
  \end{tagenumerate}
\end{proof}

\begin{probtag}{probEx,probAll}
  أكمل برهان \olref[fol][syn][sem]{prop:quant-terms}.
\end{probtag}

\end{document}
